\chapter{Model Validation and Epistemic Humility}
\label{ch:model-validation}

\epigraph{The only useful function of a statistician is to make predictions, and thus to provide a basis for action.}{W.\ Edwards Deming}

\noindent
Doxonomics makes strong claims about the causal structure of belief: that material resources shape public opinion through identifiable institutional mechanisms.
These claims must be testable, falsifiable, and honest about their limitations.
A field that studies how power corrupts knowledge production must be especially vigilant about how its own knowledge production might be corrupted---by ideological commitment, confirmation bias, unfalsifiable theorizing, or the temptation to see doxonomic causation in every correlation between money and belief.

This chapter develops the tools for validating doxonomic models---uncertainty quantification, out-of-sample prediction, sensitivity analysis, triangulation, and Bayesian model comparison---and the intellectual discipline to admit when those models fail.
It is the methodological conscience of the book.

\section{Uncertainty Quantification}
\label{sec:ch20-uncertainty}

\subsection{The uncertainty imperative}

\begin{principle}[Doxonomic Uncertainty]
\label{prin:uncertainty}
\index{uncertainty quantification}
Every doxonomic estimate---of funding flows, belief prevalence, causal effects, model parameters---must be accompanied by a quantified uncertainty.
Point estimates without uncertainty intervals are not doxonomic science; they are doxonomic assertion.
\end{principle}

The principle sounds banal---every statistics course teaches confidence intervals---but in practice, doxonomic research faces distinctive challenges in uncertainty quantification:
\begin{itemize}[nosep]
  \item funding data is incomplete (dark money, international flows, in-kind contributions), so the inputs to doxonomic models carry large measurement error,
  \item belief data comes from surveys with known biases (social desirability, framing effects), so the outputs carry systematic error,
  \item causal identification is rarely bulletproof, so structural uncertainty (``is this the right model?'') dominates statistical uncertainty (``given this model, how precise is the estimate?'').
\end{itemize}

\subsection{Frequentist and Bayesian frameworks}

\begin{definition}[Doxonomic Uncertainty Interval]
\label{def:uncertainty-interval}
\index{uncertainty interval}
A $100(1-\alpha)\%$ uncertainty interval for a doxonomic parameter $\theta$ can be constructed in two frameworks:
\begin{itemize}[nosep]
  \item \textbf{Frequentist (confidence interval)}: the procedure $[\hat{\theta}_L, \hat{\theta}_U]$ is calibrated so that, over repeated sampling, the interval contains the true $\theta$ at least $100(1-\alpha)\%$ of the time.
  The probability statement is about the procedure, not about $\theta$ lying in any particular realized interval.
  \item \textbf{Bayesian (credible interval)}: given a prior $\pi(\theta)$ and data $y$, the interval $[\theta_L, \theta_U]$ satisfies $P(\theta_L \leq \theta \leq \theta_U \mid y) \geq 1 - \alpha$ under the posterior.
  The probability statement is about $\theta$ conditional on the observed data.
\end{itemize}
Doxonomic research should specify which framework is in use.
\end{definition}

Given the typically small samples, strong prior information, and complex hierarchical structures in doxonomic studies, Bayesian methods are often more natural.
A Bayesian framework allows the analyst to incorporate prior knowledge about funding magnitudes, belief distributions, and institutional behavior while updating transparently as data arrives.

\subsection{Propagating uncertainty through the doxonomic chain}

The doxonomic causal chain $F \to I \to N \to E \to B$ involves multiple estimated quantities at each link.
Uncertainty at each link compounds:

\begin{proposition}[Uncertainty Propagation]
\label{prop:uncertainty-propagation}
\index{uncertainty propagation}
If the doxonomic effect of funding $F$ on belief $B$ is the product of link-specific effects:
\[
  \frac{\partial B}{\partial F} = \frac{\partial B}{\partial E} \cdot \frac{\partial E}{\partial N} \cdot \frac{\partial N}{\partial I} \cdot \frac{\partial I}{\partial F}
\]
and each link is estimated with independent uncertainty, then the total variance of the end-to-end effect is (by the delta method):
\[
  \operatorname{Var}\!\left(\frac{\partial B}{\partial F}\right) \approx \left(\frac{\partial B}{\partial F}\right)^2 \sum_{k} \left(\frac{\sigma_k}{\hat{\beta}_k}\right)^2
\]
where $\hat{\beta}_k$ and $\sigma_k$ are the estimate and standard error at link $k$.
The relative uncertainty of the total effect is approximately the square root of the sum of squared relative uncertainties at each link.
\end{proposition}

This means that even if each individual link is estimated with moderate precision (say, a coefficient of variation of 30\%), the end-to-end estimate may have enormous uncertainty.
For a four-link chain with 30\% relative uncertainty at each link, the total relative uncertainty is $\sqrt{4 \times 0.09} = 0.60$---a 60\% coefficient of variation, yielding a very wide confidence interval.
This is a fundamental constraint on doxonomic measurement, not a failure of technique.

\subsection{Reporting standards}

We propose minimum reporting standards for doxonomic uncertainty:

\begin{enumerate}[nosep]
  \item \textbf{Primary estimates}: point estimate $\pm$ 95\% interval, with the framework (frequentist or Bayesian) stated.
  \item \textbf{Data uncertainty}: document missing data rates for funding, belief, and exposure variables.
  \item \textbf{Model uncertainty}: report results under at least two plausible model specifications.
  \item \textbf{Identification uncertainty}: state the identifying assumptions and assess their plausibility.
  \item \textbf{Propagated uncertainty}: when combining estimates across links, report the compounded interval.
\end{enumerate}

\section{Overfitting and the Ideological Hazard}
\label{sec:ch20-overfitting}

\subsection{The problem}

\begin{definition}[Doxonomic Overfitting]
\label{def:overfitting}
\index{overfitting}
A doxonomic model \emph{overfits} if it explains observed patterns in a specific dataset (e.g., media coverage correlates with funding) but fails to predict out-of-sample---because the apparent pattern reflects noise, selection, or post hoc rationalization rather than genuine causal structure.
\end{definition}

Doxonomics is especially vulnerable to overfitting for three reasons:

\begin{enumerate}
  \item \textbf{The dependent variable is over-determined.}
  Public belief is influenced by funding, institutions, media, education, personal experience, peer effects, economic conditions, historical memory, and random events.
  Any model that includes only doxonomic variables will have low $R^2$, tempting analysts to mine the specification until something ``works.''

  \item \textbf{Funding data is incomplete and noisy.}
  Missing data on dark money, in-kind contributions, and indirect funding creates measurement error that can generate spurious patterns when combined with flexible model specifications.

  \item \textbf{The analyst has priors.}
  A researcher drawn to doxonomics is likely to believe, prior to analysis, that power shapes belief.
  This prior creates a confirmation bias that makes it easier to accept evidence consistent with doxonomic explanations and harder to accept null results.
\end{enumerate}

\subsection{The ideological hazard}

\begin{principle}[Ideological Hazard]
\label{prin:ideological-hazard}
\index{ideological hazard}
A discipline that studies how ideology distorts knowledge production must acknowledge the possibility that it is itself subject to ideological distortion.
Doxonomics is not immune to the forces it studies.
The ideological hazard is the temptation to:
\begin{itemize}[nosep]
  \item see doxonomic causation in every correlation between money and belief,
  \item attribute all belief to manipulation while ignoring genuine persuasion, rational updating, and independent thought,
  \item treat one's own beliefs as exempt from doxonomic influence,
  \item present speculative hypotheses with the confidence appropriate to well-identified causal claims.
\end{itemize}
\end{principle}

The ideological hazard does not invalidate doxonomics---the same hazard exists in any social science that studies power---but it requires institutional safeguards: pre-registration, adversarial collaboration, publication of null results, and the kind of epistemic humility developed in Section~\ref{sec:ch20-humility}.

\section{Out-of-Sample Prediction}
\label{sec:ch20-prediction}

\subsection{The prediction test}

The most powerful safeguard against overfitting is out-of-sample prediction.
A model that can predict doxonomic outcomes in data it has not seen has demonstrated that it has captured genuine structure, not noise.

\begin{definition}[Predictive Validation]
\label{def:predictive-validation}
\index{predictive validation}
Split the available data into training set $\mathcal{D}_{\text{train}}$ and test set $\mathcal{D}_{\text{test}}$ (or use $k$-fold cross-validation).
Fit the doxonomic model on $\mathcal{D}_{\text{train}}$ and evaluate on $\mathcal{D}_{\text{test}}$ using:
\begin{itemize}[nosep]
  \item \textbf{Root mean squared prediction error} (RMSPE) for continuous outcomes:
  \[
    \text{RMSPE} = \sqrt{\frac{1}{|\mathcal{D}_{\text{test}}|} \sum_{i \in \mathcal{D}_{\text{test}}} (Y_i - \hat{Y}_i)^2}
  \]
  \item \textbf{Classification accuracy} (or AUC) for binary belief outcomes.
  \item \textbf{Log predictive density} for Bayesian models: $\sum_{i \in \mathcal{D}_{\text{test}}} \log p(Y_i \mid \mathcal{D}_{\text{train}})$.
\end{itemize}
Compare the doxonomic model's prediction against a \emph{null model} that uses only non-doxonomic predictors (demographics, economic conditions, prior beliefs).
The doxonomic model must outperform the null to justify its additional complexity.
\end{definition}

\subsection{Temporal out-of-sample tests}

In doxonomic settings, temporal splits are more informative than random splits.
Train the model on data up to year $t$ and predict outcomes for years $t+1, t+2, \ldots$
This tests whether the model's structural relationships are stable over time---a stronger test than cross-sectional prediction.

\begin{example}
\label{ex:temporal-prediction}
A model relates think-tank funding to public opinion on climate change, estimated on 1990--2010 data.
The model predicts that the decline in fossil-fuel think-tank funding after 2015 should lead to a modest increase in climate concern.
Checking this prediction against 2015--2022 survey data:

\begin{center}
\begin{tabular}{lcc}
\toprule
Year & Predicted climate concern & Actual climate concern \\
\midrule
2016 & 62.3\% & 64.1\% \\
2018 & 64.7\% & 67.2\% \\
2020 & 66.1\% & 70.8\% \\
2022 & 67.4\% & 72.3\% \\
\bottomrule
\end{tabular}
\end{center}
The model correctly predicts the direction and approximate magnitude of the change but systematically underpredicts, suggesting that factors beyond think-tank funding (e.g., direct experience of extreme weather, youth activism) contributed to the belief shift.
The prediction error is informative: it reveals the limits of a funding-only doxonomic model.
(Illustrative data.)
\end{example}

\subsection{Cross-national out-of-sample tests}

An even stronger test: estimate the model on one country and predict outcomes in another.
If the doxonomic model captures generalizable mechanisms, a model estimated on US data should have some predictive power for UK or Australian outcomes (after adjusting for institutional differences).
Failure to generalize cross-nationally suggests that the model has overfit to country-specific features.

\section{Sensitivity Analysis}
\label{sec:ch20-sensitivity}

\subsection{Rosenbaum bounds revisited}

\begin{model}[Sensitivity Analysis for Doxonomic Claims]
\label{mod:sensitivity}
\index{sensitivity analysis}
Following \textcite{rosenbaum2002}, a sensitivity analysis asks: how large would an unobserved confounder $U$ have to be to explain away the estimated doxonomic effect?
The sensitivity parameter $\Gamma$ represents the maximum ratio of treatment odds for two units with the same observed covariates:
\[
  \frac{1}{\Gamma} \leq \frac{P(D=1 \mid X, U=1) / P(D=0 \mid X, U=1)}{P(D=1 \mid X, U=0) / P(D=0 \mid X, U=0)} \leq \Gamma.
\]
At $\Gamma = 1$, there is no unobserved confounding.
As $\Gamma$ increases, the bounds on the treatment effect widen.
The \emph{breaking point} $\Gamma^*$ is the smallest $\Gamma$ at which the doxonomic effect is no longer statistically significant.
A claim with $\Gamma^* = 1.2$ is fragile; one with $\Gamma^* = 3.0$ is robust.
\end{model}

\subsection{The Oster test}

\textcite{oster2019} provides an alternative sensitivity framework based on coefficient stability:

\begin{proposition}[Coefficient Stability Test]
\label{prop:oster-test}
\index{coefficient stability}
Let $\tilde{\tau}$ be the doxonomic effect estimate from a short regression (without controls), $\hat{\tau}$ the estimate from a long regression (with observable controls), and $\tilde{R}^2, \hat{R}^2$ the corresponding $R$-squared values.
Under the assumption that observable and unobservable confounders are equally important ($\delta = 1$) and that the maximum attainable $R^2$ is $R_{\max}$, the bias-adjusted estimate is:
\[
  \tau^* = \hat{\tau} - (\tilde{\tau} - \hat{\tau}) \cdot \frac{R_{\max} - \hat{R}^2}{\hat{R}^2 - \tilde{R}^2}.
\]
If $\tau^* > 0$ (or retains the hypothesized sign) for $R_{\max} = 1.3 \hat{R}^2$ and $\delta = 1$, the estimate is considered robust.
\end{proposition}

\begin{example}
\label{ex:oster-application}
A study finds that think-tank proximity predicts pro-deregulation beliefs:
\begin{itemize}[nosep]
  \item Short regression (no controls): $\tilde{\tau} = 0.42$, $\tilde{R}^2 = 0.05$.
  \item Long regression (with demographics, income, education, urban/rural): $\hat{\tau} = 0.28$, $\hat{R}^2 = 0.18$.
  \item With $R_{\max} = 1.3 \times 0.18 = 0.234$ and $\delta = 1$:
  \[
    \tau^* = 0.28 - (0.42 - 0.28) \cdot \frac{0.234 - 0.18}{0.18 - 0.05} = 0.28 - 0.14 \cdot \frac{0.054}{0.13} = 0.28 - 0.058 = 0.222.
  \]
\end{itemize}
The bias-adjusted estimate $\tau^* = 0.222$ remains positive and substantively meaningful.
Even if unobserved confounders are as important as observed ones, the doxonomic effect survives.
\end{example}

\subsection{Placebo tests}

\begin{definition}[Doxonomic Placebo Test]
\label{def:placebo}
\index{placebo test}
A \emph{placebo test} applies the doxonomic estimation procedure to an outcome that should \emph{not} be affected by the hypothesized mechanism.
If a significant effect is found on the placebo outcome, the original finding is likely driven by confounding rather than the doxonomic mechanism.
\end{definition}

\begin{example}
\label{ex:placebo}
A study claims that fossil-fuel funding shifted beliefs about climate change.
Placebo tests:
\begin{enumerate}[nosep]
  \item Does fossil-fuel funding predict beliefs about \emph{evolution}?
  If yes, the ``effect'' may reflect general conservative ideology, not a specific doxonomic mechanism.
  \item Does fossil-fuel funding predict beliefs about climate change \emph{before} the funding began?
  If yes, the relationship is driven by pre-existing differences (selection), not the funding.
  \item Does fossil-fuel funding predict climate beliefs in countries where the funded institutions have no presence?
  If yes, the correlation reflects global trends, not institutional influence.
\end{enumerate}
Passing all three placebo tests substantially strengthens the doxonomic claim.
\end{example}

\section{Bayesian Model Comparison}
\label{sec:ch20-bayesian}

\subsection{The model comparison problem}

Doxonomic claims typically involve choosing among competing explanations for the same observed pattern:
\begin{itemize}[nosep]
  \item $M_1$: funding causes belief change (the doxonomic explanation),
  \item $M_2$: ideology drives both funding and belief (confounding),
  \item $M_3$: belief change drives funding (reverse causation---funders invest where they see receptive audiences),
  \item $M_4$: no systematic relationship (coincidence).
\end{itemize}

Bayesian model comparison provides a principled framework for adjudicating among these.

\begin{definition}[Bayes Factor]
\label{def:bayes-factor}
\index{Bayes factor}
The Bayes factor comparing model $M_1$ to model $M_2$ is:
\[
  \text{BF}_{12} = \frac{P(\text{data} \mid M_1)}{P(\text{data} \mid M_2)} = \frac{\int P(\text{data} \mid \theta_1, M_1) \pi(\theta_1 \mid M_1) \, d\theta_1}{\int P(\text{data} \mid \theta_2, M_2) \pi(\theta_2 \mid M_2) \, d\theta_2}.
\]
A Bayes factor of $\text{BF}_{12} > 10$ provides ``strong'' evidence for $M_1$; $\text{BF}_{12} > 100$ provides ``decisive'' evidence (Jeffreys' scale).
The Bayes factor automatically penalizes model complexity through the prior-weighted integral, implementing a natural Occam's razor.
\end{definition}

\subsection{Comparing doxonomic and non-doxonomic models}

The critical comparison is $M_1$ (doxonomic model: funding $\to$ belief) vs.\ $M_2$ (confounding model: ideology $\to$ funding and ideology $\to$ belief):

\begin{example}
\label{ex:bayes-factor}
Using panel data on think-tank funding and regional beliefs about deregulation:

\begin{center}
\begin{tabular}{lcc}
\toprule
Model & Log marginal likelihood & Bayes factor vs.\ $M_4$ \\
\midrule
$M_1$: funding $\to$ belief & $-3{,}412.7$ & $245.3$ \\
$M_2$: ideology confounds & $-3{,}415.1$ & $22.8$ \\
$M_3$: reverse causation & $-3{,}418.9$ & $0.5$ \\
$M_4$: no relationship & $-3{,}419.6$ & $1.0$ (reference) \\
\bottomrule
\end{tabular}
\end{center}

Both $M_1$ and $M_2$ are strongly supported over no-relationship.
The Bayes factor of $M_1$ over $M_2$ is $245.3 / 22.8 = 10.8$, providing strong (but not decisive) evidence for the doxonomic over the confounding explanation.
The reverse-causation model is not supported.
(Illustrative data.)
\end{example}

\subsection{Model averaging}

When multiple models are plausible, Bayesian model averaging (BMA) provides predictions that account for model uncertainty:
\[
  P(\theta \mid \text{data}) = \sum_{m=1}^{M} P(\theta \mid \text{data}, M_m) \cdot P(M_m \mid \text{data})
\]
where $P(M_m \mid \text{data})$ is the posterior model probability.
This is especially appropriate for doxonomics because the correct model is genuinely uncertain---confounding, reverse causation, and doxonomic causation may all operate simultaneously.

\section{Triangulation}
\label{sec:ch20-triangulation}

\subsection{The convergence criterion}

\begin{principle}[Doxonomic Triangulation]
\label{prin:triangulation}
\index{triangulation}
A doxonomic claim is strengthened when it is supported by multiple independent lines of evidence.
The key categories:
\begin{enumerate}[nosep]
  \item \textbf{Accounting evidence}: money was spent (foundation records, tax filings).
  \item \textbf{Production evidence}: narratives were produced (publications, media content).
  \item \textbf{Network evidence}: institutional connections exist between funder and producer.
  \item \textbf{Causal evidence}: exposure changes beliefs (experiments, natural experiments).
  \item \textbf{Textual evidence}: narrative content shifted in the predicted direction (text analysis).
  \item \textbf{Historical evidence}: the timing of funding precedes the timing of belief change.
  \item \textbf{Comparative evidence}: similar funding produces similar effects in different contexts.
\end{enumerate}
No single line is sufficient.
The claim is credible when multiple independent lines converge on the same conclusion.
\end{principle}

\subsection{Formal triangulation assessment}

\begin{model}[Triangulation Scorecard]
\label{mod:triangulation-score}
\index{triangulation scorecard}
For a specific doxonomic claim, rate the evidence on each line as: strong ($++$), moderate ($+$), weak ($\circ$), absent ($-$), or contradictory ($--$).
The overall assessment:
\begin{itemize}[nosep]
  \item \textbf{Established}: 4+ lines at strong or moderate, none contradictory.
  \item \textbf{Supported}: 2--3 lines at strong or moderate, none contradictory.
  \item \textbf{Suggestive}: 1--2 lines at moderate, others absent.
  \item \textbf{Speculative}: only weak evidence or single line.
  \item \textbf{Challenged}: at least one contradictory line.
\end{itemize}
\end{model}

\begin{example}
\label{ex:triangulation}
Scorecard for ``fossil-fuel funding produced climate doubt'':

\begin{center}
\begin{tabular}{lcc}
\toprule
Evidence line & Rating & Key source \\
\midrule
Accounting & $++$ & Tax filings, litigation discovery \\
Production & $++$ & Think-tank publications, op-eds \\
Network & $++$ & Board overlaps, funding flows \\
Causal & $+$ & Natural experiments, some RCTs \\
Textual & $++$ & ``Uncertainty'' frame six times more prevalent \\
Historical & $++$ & Internal memos date strategy to 1989 \\
Comparative & $+$ & Weaker in countries without think-tank network \\
\bottomrule
\end{tabular}
\end{center}
Assessment: \textbf{Established}.
Six lines at strong or moderate, none contradictory.
This is one of the best-documented doxonomic cases because litigation discovery provided archival evidence that is normally inaccessible.
\end{example}

\section{Falsifiability and Degenerating Research Programs}
\label{sec:ch20-falsifiability}

\subsection{What would falsify doxonomics?}

A field that cannot specify conditions under which its core claims would be falsified is not science.
Doxonomics should state explicitly what observations would count as evidence against its hypotheses:

\begin{enumerate}
  \item \textbf{Funding-belief independence}: if careful causal studies consistently find no relationship between funding and belief (after correcting for confounding), the core doxonomic hypothesis is falsified.

  \item \textbf{Symmetry failure}: if well-funded narratives are no more prevalent than poorly funded narratives of equal quality, the power-shapes-belief mechanism is not operative.

  \item \textbf{Resilience of reason}: if populations consistently resist funded narratives and converge on true beliefs regardless of the institutional environment, the doxonomic model of belief formation is wrong.

  \item \textbf{Reverse dominance}: if belief change consistently precedes (rather than follows) funding and institutional activity, the causal arrow runs the other direction.
\end{enumerate}

None of these falsifications would be established by a single study---they would require a pattern of findings across multiple contexts and methods.

\subsection{Lakatos and degenerating programs}

Following \textcite{lakatos1978}, a research program \emph{progresses} if it successfully predicts novel facts and \emph{degenerates} if it can only accommodate observations post hoc.
Doxonomics should be evaluated by this standard:

\begin{itemize}[nosep]
  \item A \emph{progressive} sign: a doxonomic model estimated on 1990--2010 data correctly predicts the direction and approximate magnitude of belief change in 2010--2020.
  \item A \emph{degenerating} sign: every counter-example to a doxonomic claim is explained away by ad hoc auxiliary hypotheses (``the funding was dark money we can't trace,'' ``the effect operates at a level we can't measure,'' ``the alternative explanation is itself a doxonomic product'').
\end{itemize}

The danger of degeneration is especially acute for doxonomics because its subject matter---the production of belief by power---provides a built-in excuse for any failure: ``the reason you don't see the doxonomic effect is that the doxonomic system has prevented you from seeing it.''
This is the epistemic equivalent of an unfalsifiable conspiracy theory, and guarding against it requires the formal validation tools of this chapter.

\section{Epistemic Humility in Practice}
\label{sec:ch20-humility}

\subsection{The humility checklist}

\begin{principle}[Epistemic Humility]
\label{prin:humility}
\index{epistemic humility}
Every doxonomic publication should explicitly state:
\begin{enumerate}[nosep]
  \item \textbf{What it has shown}: the specific empirical finding, described in modest terms (correlation, association, estimated causal effect under stated assumptions).
  \item \textbf{What it assumes}: the identifying assumptions, measurement assumptions, and theoretical priors that underpin the analysis.
  \item \textbf{What it has not shown}: the causal claims that the evidence does not support, even if the analyst suspects they are true.
  \item \textbf{How it could be wrong}: the most plausible alternative explanation for the observed pattern.
  \item \textbf{What would change the conclusion}: the specific evidence that would cause the author to abandon or substantially revise the claim.
\end{enumerate}
Over-claiming undermines the field's credibility---its own epistemic capital.
A field that studies the corruption of epistemic authority must guard its own epistemic authority with unusual care.
\end{principle}

\subsection{Graduated confidence language}

We propose a graduated vocabulary for doxonomic claims, tied to the triangulation scorecard:

\begin{center}
\begin{tabular}{ll}
\toprule
Triangulation level & Appropriate language \\
\midrule
Established & ``The evidence strongly suggests that\ldots'' \\
Supported & ``There is substantial evidence that\ldots'' \\
Suggestive & ``Preliminary evidence is consistent with\ldots'' \\
Speculative & ``It is plausible that\ldots'' or ``We hypothesize that\ldots'' \\
Challenged & ``The evidence is mixed; some findings suggest\ldots'' \\
\bottomrule
\end{tabular}
\end{center}

This vocabulary is not merely stylistic.
It encodes real epistemic information and allows readers to calibrate their confidence appropriately.
A field-wide norm of graduated language would do more to protect doxonomics from overclaiming than any number of statistical tests.

\subsection{The symmetry obligation}

\begin{principle}[Analytical Symmetry]
\label{prin:symmetry}
\index{analytical symmetry}
Doxonomic analysis should be applied symmetrically to all belief-producing institutions, not selectively to those the researcher opposes.
If the tools developed in this book can show that fossil-fuel funding shaped climate beliefs, the same tools should be applied to environmental foundation funding, government science communication, and academic knowledge production.
The asymmetric application of doxonomic analysis---scrutinizing only one side's narrative infrastructure---is itself a form of doxonomic bias.
\end{principle}

This is not a ``both sides'' equivalence.
Different institutions may have very different levels of doxonomic power, and the analysis may well find that one side's infrastructure is vastly larger or more effective.
But the \emph{method} of analysis must be applied symmetrically; only the \emph{results} should be asymmetric.

\section{Common Pitfalls and How to Avoid Them}
\label{sec:ch20-pitfalls}

We conclude with a catalog of the most common errors in doxonomic reasoning, distilled from the methodological principles developed throughout this part of the book.

\begin{enumerate}
  \item \textbf{Post hoc ergo propter hoc.}
  ``Think tank was founded in 1985; beliefs shifted in 1990; therefore the think tank caused the shift.''
  Temporal precedence is necessary but not sufficient for causation.
  Remedy: use the causal methods of Chapter~\ref{ch:causal-inference}.

  \item \textbf{Guilt by association.}
  ``Institution $X$ receives funding from corporation $Y$; therefore $X$'s conclusions are wrong.''
  Funding creates a \emph{suspicion} of bias, not proof of it.
  The doxonomic question is whether funding predicts framing \emph{conditional on the evidence}---not whether funded conclusions are false.
  Remedy: distinguish between the genetic fallacy (dismissing claims because of their origin) and legitimate doxonomic analysis (studying how origin correlates with content).

  \item \textbf{Monocausal explanation.}
  ``Belief $p$ exists because of doxonomic campaign $X$.''
  Real beliefs have multiple causes: personal experience, peer influence, institutional exposure, economic conditions, and yes, funded narrative production.
  The doxonomic contribution is one factor among many.
  Remedy: estimate the \emph{partial} effect of doxonomic variables, controlling for other factors.

  \item \textbf{Unfalsifiable theorizing.}
  ``If you can't see the doxonomic influence, that's because the influence is so total that it has become invisible.''
  This renders the theory immune to evidence.
  Remedy: specify observable implications and test them (Section~\ref{sec:ch20-falsifiability}).

  \item \textbf{The totalizing error.}
  ``All belief is socially produced; therefore no belief is genuinely held.''
  Showing that a belief has social origins does not show that it is false or that the believer is a puppet.
  People are agents who process, filter, resist, and transform the narratives they encounter.
  Remedy: model agent heterogeneity in belief response (Chapter~\ref{ch:formal-models}).

  \item \textbf{Asymmetric scrutiny.}
  Applying doxonomic analysis only to beliefs one disagrees with.
  Remedy: the symmetry obligation (Principle~\ref{prin:symmetry}).
\end{enumerate}

\section{Notes and References}
\label{sec:ch20-notes}

Model selection and validation are covered in \textcite{freedman2009} and \textcite{gelman2020}.
Sensitivity analysis via Rosenbaum bounds is developed in \textcite{rosenbaum2002}; the Oster test is in \textcite{oster2019}.
On the limits of causal inference in social science, see \textcite{manski1993} and \textcite{imbens2015}.

Bayesian model comparison and Bayes factors are treated in \textcite{kass1995} and Gelman et al.\ (2013).
Model averaging is developed in Hoeting et al.\ (1999).
The philosophy of falsifiability and research programs draws on \textcite{lakatos1978} and Popper (1963).

For triangulation methodology, see \textcite{king1994} and \textcite{morgan2015}.
On the replication crisis and its implications for social science methodology, see Open Science Collaboration (2015).
The ethics of social science research are discussed in \textcite{salganik2018}.

On the problem of reflexivity in social science---the fact that studying social phenomena can change them---see \textcite{hacking1999} and Bourdieu (2004).

\section{Exercises}

\begin{exercise}
Take a published claim of the form ``think-tank funding increased public support for deregulation.''
\begin{enumerate}[nosep]
  \item Design a Rosenbaum sensitivity analysis. For what values of $\Gamma$ does the claim survive?
  \item Apply the Oster coefficient stability test. Compute $\tau^*$ for $R_{\max} = 1.3\hat{R}^2$ and $\delta = 1$.
  \item Design three placebo tests that would strengthen or weaken the claim.
\end{enumerate}
\end{exercise}

\begin{exercise}
Describe a doxonomic hypothesis that is unfalsifiable in practice.
\begin{enumerate}[nosep]
  \item Explain why it is unfalsifiable.
  \item Reformulate it as a testable hypothesis with specific observable implications.
  \item Identify the data and methods needed to test the reformulated hypothesis.
  \item Describe what results would constitute falsification.
\end{enumerate}
\end{exercise}

\begin{exercise}
Construct a triangulation scorecard (Model~\ref{mod:triangulation-score}) for the claim that economics education increases belief in selfish human nature.
Rate each of the seven evidence lines and provide an overall assessment.
Where is the evidence weakest?
What study would most improve the overall assessment?
\end{exercise}

\begin{exercise}
\label{ex:bayes-factor-exercise}
Consider two models for the correlation between media ownership concentration and political polarization:
\begin{itemize}[nosep]
  \item $M_1$: concentration causes polarization (doxonomic model).
  \item $M_2$: polarization causes audiences to self-sort, leading to concentrated viewership that \emph{looks like} media power but is actually demand-driven.
\end{itemize}
\begin{enumerate}[nosep]
  \item What data would discriminate between these models?
  \item If you had panel data on both ownership changes and polarization measures, how would you compute the Bayes factor for $M_1$ vs.\ $M_2$?
  \item What prior over model parameters would you use, and how sensitive is the Bayes factor to this choice?
\end{enumerate}
\end{exercise}

\begin{exercise}
Apply the common pitfalls catalog (Section~\ref{sec:ch20-pitfalls}) to a real-world doxonomic claim from public discourse (e.g., ``Big Pharma controls what doctors believe,'' ``the liberal media brainwashes viewers,'' ``tech companies radicalize users'').
\begin{enumerate}[nosep]
  \item Which pitfalls does the claim commit?
  \item How would you reformulate the claim to avoid those pitfalls?
  \item What evidence would you need to evaluate the reformulated claim?
\end{enumerate}
\end{exercise}

\begin{exercise}
Design a pre-registration protocol for a doxonomic study that is maximally protected against the ideological hazard (Principle~\ref{prin:ideological-hazard}).
Include:
\begin{enumerate}[nosep]
  \item specific hypotheses with predicted directions and magnitudes,
  \item primary and secondary outcomes,
  \item the analysis plan (including what constitutes a null result),
  \item an adversarial critique section: write the strongest objection a skeptic could raise, and describe how the study design addresses it.
\end{enumerate}
\end{exercise}

\begin{exercise}[Computational]
Simulate the uncertainty propagation problem (Proposition~\ref{prop:uncertainty-propagation}):
\begin{enumerate}[nosep]
  \item Create a four-link doxonomic chain where each link has a true effect of $\beta = 0.5$ and is estimated with noise $\varepsilon \sim N(0, 0.15^2)$.
  \item Compute the end-to-end effect $\hat{\beta}_{\text{total}} = \prod_{k=1}^{4} \hat{\beta}_k$ for 10,000 Monte Carlo replications.
  \item Plot the sampling distribution of $\hat{\beta}_{\text{total}}$.
  \item What fraction of replications yield an estimate within 50\% of the true total effect ($0.5^4 = 0.0625$)?
  \item How does this fraction change if you increase the per-link noise to $\sigma = 0.25$?
\end{enumerate}
\end{exercise}
